\section{前言}
截止2026年9月16日,openAI已经攻克多个数学难题,其中甚至包括千禧年难题,那么在当前这个环境的冲击下,
不得不反思我们所做的工作的意义,诚然,本次实验利用gpt6就能很快作出.现在事实上进入了一个产出远远大于消化的时代,
我们难道就只是做一个在电脑前yes or no的机器人?

从历史上说,一开始的编程语言就是打孔机,然后变成了机器语言,汇编语言,最后才到了如今的编程语言,每个时代都有每个时代的古法编程,
总不至于如今不学C语言,而是不断地写0101?

但是AI似乎和这个又不太一样,即使过去有各种各样的工具提升,但是像AI这样“自举”的感觉是没有的,举例来说,C语言本身不能教你学会C语言,
但是如果你不会使用AI,AI可以教你使用AI。也就是AI变成了一个类似许愿机的万能工具,而且是可以用自然语言驱动的,这是史无前例的

那么也就是说,实验报告可以AI写,实验代码可以AI生成.从我个人看来,我已经完全接受AI进行排版之类的琐事,比如latex代码的编写,这个显然对于
现在是没有什么意义的,但是完全让AI生成内容,可能接受度还不是很高。因为如果让AI完全生成内容的话,那这样AI就完全变成了一个许愿池,变成了一场沙里淘金的游戏

也就是变成了一种运气的比拼,我们所要做的只是从浩如烟海的概率生成的AI文本中挖掘出有用的部分。但是这样我又有另一个想法,那‘挖掘’这项工作又真的需要人来做吗？
我们完全可以把一个AI生成的内容发给另一个AI,让他判断其中有价值的部分,也就是agent合作,openAI在解决ns问题时就使用了10000个agent,这其中并没有人类的参与

一个残酷的事实是,在人与AI交互的过程中,最拖累项目进程的或许就是人,如果把人换成别的AI,效率就会显著提升

不过还有一种观点认为,AI不能进行价值判断,到最后还是要人来审阅价值,也就是人被异化成一个打分的工具,但是我们对AI的反馈似乎也是被作为AI训练的数据之一,人的判断本身
也在AI学习的范围内,这也在前面说的agent合作体现了。所以这个价值不应该是客观的价值,而应该是对人的价值。诚然,AI可以产出很多有客观价值的东西,但是这些客观价值是否有人的主观上的价值,这就真的需要人来了。

回归到实验报告本身,AI固然是可以生成一份看上去完美无暇的报告,并且我预计在本次实验中,大多数报告看上去会是这样的,大家都用AI,那生成的报告的差异性是比较小的。
那我们不禁反思这样完美无瑕的报告真的有意义吗?有人可能说工作中哪有那么多意义？一份完美的报告能应付领导,搞好形式主义就行。但是这样难免又违背了本心,也就是结果论和过程论的冲突

反思之际,我又萌生了一个想法,既然在AI时代,答案都是一样的完美正确,但是过程呢？

值得注意的是,我在日常使用AI的过程中,发现了一个现象,AI生成的文本是几乎不会碰壁的,具体来说,真理就像走迷宫,除非你有上帝视角,不然不可能一次找到通途

换而言之,碰壁几乎是百分百的事情,但是AI会把碰壁藏在思维链当中(他最后生成的文本几乎是不会告诉你,他在解决问题的过程中遇到了什么困难,哪条路行不通,当然DeepSeek的思维链是可以查看的,但是我们往往不会关注)

回到前面说的,既然正确的答案是轻而易举的,那么错误的答案自然就有了一些价值,毕竟物以稀为贵嘛,我们人本身是不可能模仿AI的思考方式的,所以我们一定会碰壁

那么既然如此,AI生成的文本某种程度上来说,对于人类的思维是不可复现的,那份完美无瑕的报告,人类的思维很难做到

那我们能做到什么？就是说明自己的碰的壁,也就是错误,令人感到反差的事,错误居然反而更有价值,因为错误才能让人反思,那这份实验报告才有学习意义

从我个人而言,目前的观点是,在学习的时候应该将AI当作老师和搜索引擎,而在工程实践中,你确实可以把它当做你的员工来看待,但是必须review,否则是没有意义的

所以我认为现在的本次实验报告,应该体现我思考问题的时候,哪里碰壁了,为什么会碰壁,最后又是如何找到正确的思维路径的(这当然可以通过AI来纠正),包括对AI的质疑(不过对于如今的模型AI很难出错)

总而言之,就是要把出错到纠正的这条心路历程写出来,这样的文章对人类来说是有意义的,因为这样的思维链路可能是大多数人都会遇到的,是比较有普适性的

另一方面,值得注意的一件事是,我在学习的过程中(包括大学三年和以往所有经历),发现了一件有趣的事,人类其实是不太擅长处理文本的,往往一张图可以胜过千言万语,这就是ppt等展示的意义,包括生活中肢体语言有时候比嘴上的话更容易让人明白意思

数学角度的例子就是,如果给一个函数图像(有限长度的),我们往往可以一眼从全局视角看到函数的最低点,但是计算机是不行的,他可能要通过求导解方程,梯度下降法等等算法来找到这个最低点

不过人类也不知道这个最低点到底等于多少,只是感性地察觉到“他就在那里”,这种感觉往往也是很重要的

基于这样的哲学,本次实验报告会利用gpt-image2.0多画一些示意图,并且会解释我在学习过程中过遇到的一些误区和迷思